% !TEX root = ../../main.tex

\section{系统总体概述}

系统需求分析是系统设计与实现的基础，其主要任务是明确系统面向的用户角色、核心功能范围、主要业务流程以及非功能性要求。
本文所设计的 HearBridge 系统面向听障辅助训练与手语识别场景，目标是为普通用户提供手语资源浏览、动作示范、实时手势训练、句子视频识别和识别结果反馈等能力，同时为系统管理员提供手势资源维护、样本管理、模型训练和模型版本发布等管理能力。

从使用角色来看，系统主要包括普通用户和系统管理员两类角色。
普通用户通过 HarmonyOS 移动端使用系统，主要关注手语资源学习、训练过程体验和识别结果展示；
系统管理员通过 Web 管理端维护系统资源和识别模型，主要关注手势分类、手势资源、样本数据、训练任务和模型版本的管理。
除上述两类直接用户外，Python 手势识别服务和 DeepSeek 语义修正服务作为系统支撑能力参与业务流程，但它们不属于直接操作系统的用户角色，而是作为后续系统设计与实现章节中的服务组件进行说明。

HearBridge 系统当前的核心业务可概括为两条主线：一是面向普通用户的手语训练与识别主线，用户在移动端浏览手语资源、观看动作示范并进行实时手势识别训练；
二是面向管理员的资源与模型维护主线，管理员通过后台管理端维护手势分类、手势资源、训练样本和模型版本，从而为移动端识别功能提供数据和模型支撑。
在此基础上，系统还提供句子视频识别功能，允许用户选择本地手语视频并查看原始识别结果和语义修正结果，作为受限场景下辅助交流能力的探索。
